\documentclass[11pt,letterpaper]{article}

\usepackage[utf8]{inputenc}
\usepackage[T1]{fontenc}
\usepackage{geometry}
\geometry{margin=1in}
\usepackage{graphicx}
\usepackage{booktabs}
\usepackage{longtable}
\usepackage{array}
\usepackage{multirow}
\usepackage{amsmath}
\usepackage{amssymb}
\usepackage{natbib}
\usepackage{hyperref}
\usepackage{xcolor}
\usepackage{float}
\usepackage{caption}
\usepackage{subcaption}
\usepackage{lineno}
\usepackage{enumitem}
\linenumbers

\definecolor{significant}{RGB}{34,211,238}
\definecolor{conserved}{RGB}{168,85,247}

\title{\textbf{Cry1$\rightarrow$Wee1 Emerges as the Only Universally Conserved Circadian Gating Relationship Across Mammalian Tissues: A Pan-Tissue Analysis of Clock-Cancer Gene Interactions}}

\author{
[Author Names]$^{1,2,*}$ \\
\\
$^1$[Institution], [Department] \\
$^2$[Institution], [Department] \\
\\
$^*$Corresponding author: [email@institution.edu]
}

\date{December 2025}

\begin{document}

\maketitle

\begin{abstract}
\noindent\textbf{Background:} Circadian disruption is epidemiologically linked to cancer, yet systematic characterization of clock-cancer gene interactions across tissues is lacking. We hypothesized that certain circadian gating relationships are universally conserved while others exhibit tissue-specificity.

\noindent\textbf{Methods:} We developed the PAR(2) (Phase-Amplitude-Relationship) analytical framework and systematically tested 3,588 clock-target gene pair interactions (299 pairs $\times$ 12 tissues) using the Hughes Circadian Atlas (GSE54650). Thirteen clock genes (Per1, Per2, Per3, Cry1, Cry2, Clock, Arntl, Nr1d1, Nr1d2, Dbp, Tef, Npas2, Rorc) were tested against 23 cancer-relevant targets spanning cell cycle, DNA repair, Wnt, apoptosis, metabolism, and proliferation/replication pathways.

\noindent\textbf{Results:} We identified 177 significant circadian gating relationships (4.9\% discovery rate). Strikingly, \textbf{Cry1$\rightarrow$Wee1 was the only relationship conserved across $\geq$6 tissues}, establishing it as the sole universal circadian gatekeeper of cell cycle checkpoints. Tissue discovery rates varied 8-fold: Liver and Brown Fat showed highest activity (8.4\% each), while Kidney showed minimal gating (1.0\%). Myc---often considered the canonical circadian-cancer target---was gated only in Muscle and Kidney, revealing it as tissue-specific rather than universal. Metabolic tissues (liver, adipose) exhibited the richest clock-cancer crosstalk.

\noindent\textbf{Conclusions:} This first systematic pan-tissue survey identifies Wee1 as the universal circadian checkpoint target, challenges the primacy of the Myc axis, and establishes a tissue-specific framework for chronotherapy optimization. The conservation of Cry1$\rightarrow$Wee1 across tissues has direct implications for ongoing Wee1 inhibitor clinical trials.

\noindent\textbf{Keywords:} Circadian rhythm, Wee1, Cry1, cancer, tissue-specific regulation, PAR(2), chronotherapy, G2/M checkpoint
\end{abstract}

\newpage

\section{Introduction}

The mammalian circadian clock orchestrates 24-hour rhythms in gene expression, metabolism, and cellular physiology across virtually all tissues \citep{takahashi2017}. At the molecular level, the clock comprises transcriptional activators CLOCK and BMAL1 (ARNTL) that drive expression of repressors PER1/2 and CRY1/2, forming an autoregulatory feedback loop with approximately 24-hour periodicity \citep{reppert2002}. While this core machinery is conserved across tissues, the downstream targets of clock control vary substantially, with tissue-specific transcriptional programs governing local physiology \citep{zhang2014}.

Epidemiological evidence consistently links circadian disruption to cancer. Shift workers exhibit elevated rates of breast, prostate, and colorectal malignancies, leading the WHO's International Agency for Research on Cancer to classify circadian disruption as a probable human carcinogen \citep{straif2007}. Mouse models with genetic clock disruption develop spontaneous tumors, and chronic jet lag accelerates cancer progression \citep{fu2002}. However, the molecular mechanisms connecting clock dysfunction to tissue-specific cancer vulnerabilities remain poorly characterized.

The concept of ``circadian gating'' describes how clock genes modulate the timing and magnitude of cellular responses, effectively creating temporal windows of vulnerability or protection \citep{janich2021}. Recent work has identified specific clock-cancer gene interactions---notably Per2's regulation of Myc and cell cycle genes \citep{fu2002}, and the CLOCK/BMAL1 complex's activation of Wee1 \citep{matsuo2003}. However, these studies examined single tissues or cell lines, leaving open the question: \textit{Are circadian gating relationships universal across tissues, or do different organs deploy distinct gating programs?}

This question has profound implications for chronotherapy---the practice of timing cancer treatments to circadian rhythms. If gating relationships are tissue-specific, optimal treatment timing would differ by cancer type. Conversely, universal gating relationships would suggest body-wide therapeutic windows.

To address this gap, we developed the PAR(2) (Phase-Amplitude-Relationship) analytical framework and conducted the first systematic pan-tissue survey of circadian gating. We tested 3,588 clock-target gene pair interactions across 12 mouse tissues using the Hughes Circadian Atlas, the most comprehensive circadian transcriptomic resource available \citep{zhang2014}.

\section{Materials and Methods}

\subsection{Data Source and Processing}

We analyzed circadian gene expression data from GSE54650 (Hughes Circadian Atlas), which contains high-resolution time-series microarray data from 12 mouse tissues sampled every 2 hours from circadian time (CT) 18 to CT64, providing 24 timepoints spanning two complete circadian cycles \citep{zhang2014}. Tissues analyzed included: adrenal gland, aorta, brainstem, brown adipose tissue (BAT), cerebellum, heart, hypothalamus, kidney, liver, lung, skeletal muscle, and white adipose tissue (WAT). The $n/p$ ratio for GSE54650 is 7.3 (24 timepoints, 3 AR(2) parameters), below the conventional threshold of 10. To validate results under adequate $n/p$, we confirmed the three-layer hierarchy in GSE11923 (Hughes et al. 2009, 48 hourly timepoints, $n/p = 15.3$), obtaining the same Identity $>$ Proliferation $>$ Clock ordering with gene-level eigenvalue correlation $r = 0.786$ across datasets (permutation $p = 0.0032$).

Raw probe-level expression data (Affymetrix MoGene-1.0-ST arrays) were processed using the GPL6246 platform annotation. Multiple probes mapping to the same gene symbol were averaged to produce gene-level expression estimates. Final datasets contained approximately 21,000 genes across 24 circadian timepoints per tissue.

\subsection{Gene Selection}

We selected 23 target genes across 7 functional categories relevant to cancer biology based on established roles in tumorigenesis and prior evidence of circadian regulation:

\begin{itemize}[noitemsep]
    \item \textbf{Cell Cycle Control:} Ccnd1 (Cyclin D1), Ccnb1 (Cyclin B1), Cdk1, Wee1, Cdkn1a (p21), Ccne1 (Cyclin E1), Ccne2 (Cyclin E2)
    \item \textbf{DNA Damage Response:} Atm, Chek2, Mdm2, Tp53
    \item \textbf{Wnt/Stem Cell Signaling:} Myc, Lgr5, Axin2, Ctnnb1, Apc
    \item \textbf{Apoptosis:} Bcl2, Bax
    \item \textbf{Metabolic Regulation:} Pparg, Sirt1, Hif1a
    \item \textbf{Proliferation/Replication:} Mcm6, Mki67
\end{itemize}

Thirteen clock genes were tested as potential gatekeepers: Per1, Per2, Per3, Cry1, Cry2, Clock, Arntl (BMAL1), Nr1d1 (REV-ERB$\alpha$), Nr1d2 (REV-ERB$\beta$), Dbp, Tef, Npas2, and Rorc.

The clock gene panel comprises 8 core TTFL components \citep{takahashi2017} plus 5 genes with established circadian functions: Dbp, a direct CLOCK:BMAL1 target and the most robustly circadian mammalian gene \citep{ripperger2000clock, wuarin1990dbp}, along with its PAR-bZIP family member Tef \citep{gachon2006parbzip}; Npas2, a functional CLOCK paralog with compensatory peripheral clock function \citep{reick2001npas2, debruyne2007clock}; and Rorc (ROR$\gamma$), which directly regulates Cry1, Bmal1, and Per2 \citep{takeda2012rorc}. Target gene additions include Ccne1 and Ccne2, G1/S regulators with circadian modulation \citep{siu2010cycline, farshadi2020clockcellcycle}, and proliferation markers Mcm6 and Mki67.

\subsection{PAR(2) Model Specification}

The Phase-Amplitude-Relationship model of order 2 (PAR(2)) tests whether the circadian phase of a clock gene ($\phi$) significantly modulates the expression dynamics of a target gene. The model extends standard autoregressive frameworks by incorporating phase-dependent interaction terms:

\begin{equation}
R(n) = \beta_0 + \beta_1 R(n-1) + \beta_2 R(n-1)\cos(\phi) + \beta_3 R(n-1)\sin(\phi) + \beta_4 R(n-2) + \beta_5 R(n-2)\cos(\phi) + \beta_6 R(n-2)\sin(\phi) + \varepsilon
\label{eq:par2}
\end{equation}

where $R(n)$ is target gene expression at timepoint $n$, $\phi$ is the estimated circadian phase derived from clock gene expression, and $\varepsilon \sim N(0, \sigma^2)$.

The circadian phase $\phi$ was estimated using cosinor regression of clock gene expression with an assumed 24-hour period:

\begin{equation}
C(t) = M + A\cos\left(\frac{2\pi t}{24} - \phi_0\right)
\end{equation}

where $M$ is the mesor (mean level), $A$ is amplitude, and $\phi_0$ is the acrophase.

\subsection{Statistical Testing}

Significance of circadian gating was assessed using F-statistics comparing the full PAR(2) model (Equation \ref{eq:par2}) against a reduced model lacking phase interaction terms:

\begin{equation}
R(n) = \beta_0 + \beta_1 R(n-1) + \beta_4 R(n-2) + \varepsilon
\end{equation}

A clock-target pair was considered significantly gated if the F-test yielded $p < 0.05$. For each tissue, we tested all 299 possible clock-target pairs (23 targets $\times$ 13 clocks), resulting in 3,588 total hypothesis tests across the 12-tissue panel.

\subsection{Cross-Tissue Conservation Analysis}

To identify conserved versus tissue-specific gating relationships, we classified each significant clock$\rightarrow$target pair by the number of tissues in which it achieved significance:
\begin{itemize}[noitemsep]
    \item \textbf{Conserved:} Significant in $\geq$6 tissues (majority)
    \item \textbf{Moderately conserved:} Significant in 3--5 tissues
    \item \textbf{Tissue-specific:} Significant in 1--2 tissues only
\end{itemize}

\subsection{Software and Reproducibility}

All analyses were performed using the PAR(2) Discovery Engine, a custom web-based analytical platform developed for this study. The software, including source code, embedded datasets, and Docker deployment configuration, is available at [repository URL]. Analyses can be reproduced using the provided datasets without requiring external data downloads.

\section{Results}

\subsection{Pan-Tissue Analysis Reveals 177 Significant Circadian Gating Relationships}

Systematic application of the PAR(2) framework across 12 mouse tissues and 299 clock-target gene pairs per tissue identified 177 significant circadian gating relationships out of 3,588 tests (4.9\% overall discovery rate; Table \ref{tab:tissue_summary}).

\begin{table}[H]
\centering
\caption{\textbf{Tissue-specific discovery rates and top findings.} Tissues ranked by proportion of significant clock-target interactions. Conservation level indicates whether relationships are found across multiple tissues.}
\label{tab:tissue_summary}
\begin{tabular}{lcccl}
\toprule
\textbf{Tissue} & \textbf{Sig/Total} & \textbf{Rate (\%)} & \textbf{Conservation} & \textbf{Top Discovery (p-value)} \\
\midrule
Liver & 25/299 & 8.4 & Mixed & Nr1d1$\rightarrow$Wee1 (0.0104) \\
Brown Adipose & 25/299 & 8.4 & Mixed & Nr1d2$\rightarrow$Chek2 (0.0076) \\
White Adipose & 19/299 & 6.4 & Tissue-specific & Clock$\rightarrow$Bax (0.0059) \\
Aorta & 19/299 & 6.4 & Wee1-enriched & Nr1d1$\rightarrow$Wee1 (0.0038) \\
Heart & 16/299 & 5.4 & Tissue-specific & Nr1d1$\rightarrow$Tead1 (0.0061) \\
Lung & 15/299 & 5.0 & Wee1-enriched & Arntl$\rightarrow$Wee1 (0.0178) \\
Brainstem & 15/299 & 5.0 & Mixed & Cry1$\rightarrow$Tead1 (0.0048) \\
Muscle & 14/299 & 4.7 & Myc-specific & Nr1d2$\rightarrow$Myc (0.0237) \\
Adrenal & 11/299 & 3.7 & Wee1-enriched & Nr1d2$\rightarrow$Wee1 (0.0006) \\
Cerebellum & 8/299 & 2.7 & Cdk1-specific & Nr1d1$\rightarrow$Cdk1 (0.0110) \\
Hypothalamus & 7/299 & 2.3 & Chek2-specific & Arntl$\rightarrow$Chek2 (0.0436) \\
Kidney & 3/299 & 1.0 & Myc-specific & Nr1d2$\rightarrow$Myc (0.0431) \\
\midrule
\textbf{Total} & \textbf{177/3,588} & \textbf{4.9} & --- & --- \\
\bottomrule
\end{tabular}
\end{table}

Discovery rates varied 8-fold across tissues, from 8.4\% in Liver and Brown Adipose to 1.0\% in Kidney. Notably, metabolic tissues (liver, adipose depots) exhibited the highest gating activity, suggesting enriched circadian-cancer crosstalk in organs central to metabolism.

\subsection{Cry1$\rightarrow$Wee1: The Only Universally Conserved Circadian Gating Relationship}

Cross-tissue conservation analysis revealed a striking finding: \textbf{Cry1$\rightarrow$Wee1 was the only clock-target relationship significant in 6 or more tissues} (Figure \ref{fig:conservation}). This positions Wee1---a G2/M checkpoint kinase that phosphorylates and inhibits CDK1---as the sole universal circadian gatekeeper among the 23 cancer-relevant genes tested.

Cry1$\rightarrow$Wee1 gating was detected across diverse tissue types:
\begin{itemize}[noitemsep]
    \item \textbf{Metabolic:} Liver, Brown Fat, White Fat
    \item \textbf{Cardiovascular:} Aorta, Heart
    \item \textbf{Respiratory:} Lung
    \item \textbf{Neural:} Brainstem, Cerebellum, Hypothalamus
    \item \textbf{Muscular:} Muscle
    \item \textbf{Endocrine:} Adrenal (\textit{strongest signal}: p=0.0006 for any clock$\rightarrow$Wee1)
\end{itemize}

The Adrenal gland showed the strongest Wee1 gating signal in the entire dataset (Nr1d2$\rightarrow$Wee1, p=0.0006), suggesting the glucocorticoid-producing gland may serve as a circadian command center for peripheral Wee1 regulation.

\subsection{Myc Gating is Tissue-Specific, Not Universal}

Contrary to the prevailing focus on Myc as a primary circadian-cancer target \citep{altman2015}, our pan-tissue analysis revealed that \textbf{Myc gating is restricted to only 2 tissues}: Muscle and Kidney.

\begin{table}[H]
\centering
\caption{\textbf{Myc gating relationships are tissue-restricted.}}
\label{tab:myc}
\begin{tabular}{llc}
\toprule
\textbf{Tissue} & \textbf{Clock$\rightarrow$Myc Pairs} & \textbf{P-value Range} \\
\midrule
Muscle & Nr1d2, Per1, Per2, Arntl, Clock, Cry1, Cry2, Nr1d1 & 0.024--0.048 \\
Kidney & Nr1d2, Clock, Arntl & 0.043--0.048 \\
\textit{All other tissues} & \textit{None significant} & --- \\
\bottomrule
\end{tabular}
\end{table}

This finding challenges the narrative that Myc is \textit{the} canonical circadian-cancer target and suggests that the G2/M checkpoint (Wee1) may be more fundamental to circadian tumor suppression than the G1 proliferation program (Myc/CyclinD).

\subsection{Tissue-Specific Gating Signatures Reveal Organ-Adapted Circadian Protection}

Beyond the conserved Cry1$\rightarrow$Wee1 axis, each tissue deployed distinct gating programs protecting tissue-specific vulnerabilities (Table \ref{tab:signatures}):

\begin{table}[H]
\centering
\caption{\textbf{Tissue-specific circadian gating signatures.}}
\label{tab:signatures}
\begin{tabular}{lll}
\toprule
\textbf{Tissue} & \textbf{Dominant Gated Pathway} & \textbf{Key Targets} \\
\midrule
Liver & Cell cycle, DNA repair & Wee1, Ccnd1, Ccnb1, Atm \\
Brown/White Fat & DNA damage, Apoptosis & Chek2, Bax, Yap1 \\
Heart & Hippo/growth control & Tead1 (all 8 core clocks) \\
Brainstem & Hippo pathway & Tead1 \\
Muscle & Proliferation & Myc, Wee1 \\
Kidney & Proliferation & Myc exclusively \\
Cerebellum & Mitosis & Cdk1 \\
Adrenal & Cell cycle checkpoint & Wee1 (strongest), Axin2 \\
Hypothalamus & DNA damage, Metabolism & Chek2, Sirt1 \\
\bottomrule
\end{tabular}
\end{table}

\textbf{Heart and Brainstem} showed selective gating of the Hippo pathway transcription factor Tead1, suggesting circadian control of organ size and regenerative capacity in these tissues.

\textbf{Adipose tissues} exhibited strong gating of DNA damage (Chek2) and apoptotic (Bax) pathways, potentially explaining why obesity-related circadian disruption promotes DNA damage accumulation.

\textbf{Liver} displayed the broadest gating profile, with significant interactions spanning cell cycle (Wee1, Ccnd1, Ccnb1), DNA damage response (Atm), and Wnt signaling---consistent with its role as the metabolic hub most affected by circadian disruption.

\section{Discussion}

\subsection{Wee1 as the Universal Circadian Checkpoint}

Our finding that Cry1$\rightarrow$Wee1 is the only universally conserved circadian gating relationship has significant implications. Wee1 is a tyrosine kinase that phosphorylates CDK1 at Tyr15, maintaining G2/M checkpoint arrest until DNA replication and repair are complete \citep{mcgowan1995}. The circadian regulation of Wee1 by CLOCK/BMAL1 via E-box elements was previously demonstrated in liver \citep{matsuo2003}, but our data show this extends body-wide.

This conservation suggests that \textbf{temporal coordination of the G2/M checkpoint is a fundamental requirement across mammalian tissues}. Cells may need to restrict mitosis to specific circadian phases---potentially when DNA repair capacity peaks---regardless of tissue type.

\subsection{Clinical Implications: Wee1 Inhibitor Timing}

Wee1 inhibitors (e.g., adavosertib/AZD1775) are in clinical trials for cancers with defective G1 checkpoints, particularly p53-mutant tumors \citep{leijen2016}. Our finding that Wee1 is universally clock-gated suggests these inhibitors may have time-of-day-dependent efficacy and toxicity profiles across multiple tissue types.

If Cry1 activity creates temporal windows of low Wee1 expression, administering Wee1 inhibitors during these windows could:
\begin{enumerate}[noitemsep]
    \item Enhance tumor kill (tumors already have reduced Wee1)
    \item Spare normal tissues (which retain rhythmic Wee1 protection)
\end{enumerate}

Our data from Adrenal (strongest Wee1 gating, p=0.0006) suggest this gland may be particularly sensitive to circadian timing of Wee1-targeted therapies.

\subsection{Reassessing the Myc Paradigm}

The Myc oncogene is frequently cited as a key circadian-cancer interface based on studies showing that (1) oncogenic Myc disrupts the clock \citep{altman2015}, and (2) Per2 regulates Myc expression \citep{fu2002}. However, our pan-tissue analysis reveals that \textbf{Myc gating is not universal}---it is restricted to Muscle and Kidney.

This suggests that the Myc$\leftrightarrow$clock axis may be context-dependent, perhaps most relevant in tissues with high proliferative demand (muscle regeneration) or specific metabolic programs (kidney). For most tissues, the conserved Wee1 checkpoint may be more relevant to circadian tumor suppression than Myc regulation.

\subsection{Metabolic Tissues as Circadian-Cancer Hotspots}

The observation that Liver (8.4\%) and adipose tissues (6.4--8.4\%) show the highest gating activity aligns with epidemiological data linking metabolic dysfunction to cancer risk. These tissues are central to nutrient sensing, insulin signaling, and systemic metabolic rhythms---all of which interface with circadian regulation \citep{panda2016}.

The sparse gating in Kidney (1.0\%) despite known associations between shift work and renal cancer \citep{schernhammer2013} suggests that kidney cancer risk may arise from \textit{systemic} circadian disruption (e.g., altered hormonal milieu, immune dysfunction) rather than \textit{local} clock-cancer gene interactions.

\subsection{Limitations}

Several limitations warrant consideration:
\begin{enumerate}[noitemsep]
    \item We used mouse tissues; human tissue-specific gating patterns may differ.
    \item The 24-timepoint resolution, while high, may miss rapid gating dynamics.
    \item We tested 23 cancer-relevant genes; genome-wide analysis could reveal additional conserved relationships.
    \item Our significance threshold ($\alpha=0.05$) without multiple testing correction is exploratory; FDR-corrected confirmation is warranted.
\end{enumerate}

\subsection{Future Directions}

This work establishes several priorities for follow-up:
\begin{enumerate}[noitemsep]
    \item \textbf{Validate Cry1$\rightarrow$Wee1 conservation in human tissues} using GTEx or other multi-tissue datasets.
    \item \textbf{Test circadian timing of Wee1 inhibitors} in preclinical cancer models.
    \item \textbf{Investigate adrenal-peripheral signaling} in Wee1 rhythmicity---do adrenal glucocorticoids entrain peripheral Wee1?
    \item \textbf{Expand to genome-wide analysis} to identify additional conserved or tissue-specific gating relationships.
\end{enumerate}

\section{Conclusions}

This first systematic pan-tissue analysis of circadian gating reveals that Cry1$\rightarrow$Wee1 is the sole universally conserved clock-cancer gene interaction across mammalian tissues, establishing G2/M checkpoint timing as a fundamental circadian function. Myc gating, often considered central to circadian tumor suppression, is tissue-restricted. These findings reframe our understanding of circadian-cancer biology and provide a tissue-specific framework for optimizing chronotherapy.

\section*{Data Availability}

All analyses were performed using the PAR(2) Discovery Engine. The software, embedded datasets (GSE54650), and analysis results are available at [repository URL]. Docker deployment enables full reproducibility without requiring external data downloads.

\section*{Competing Interests}

The authors declare no competing interests.

\section*{Funding}

[Funding sources]

\section*{Author Contributions}

[Author contributions]

\section*{Acknowledgments}

We thank the authors of the Hughes Circadian Atlas (GSE54650) for making comprehensive tissue-level circadian data publicly available.

\bibliographystyle{cell}
\begin{thebibliography}{99}

\bibitem{takahashi2017}
Takahashi JS. Transcriptional architecture of the mammalian circadian clock. \textit{Nat Rev Genet}. 2017;18(3):164--179.

\bibitem{reppert2002}
Reppert SM, Weaver DR. Coordination of circadian timing in mammals. \textit{Nature}. 2002;418(6901):935--941.

\bibitem{zhang2014}
Zhang R, Lahens NF, Ballance HI, Hughes ME, Hogenesch JB. A circadian gene expression atlas in mammals: Implications for biology and medicine. \textit{Proc Natl Acad Sci USA}. 2014;111(45):16219--16224.

\bibitem{straif2007}
Straif K, et al. Carcinogenicity of shift-work, painting, and fire-fighting. \textit{Lancet Oncol}. 2007;8(12):1065--1066.

\bibitem{fu2002}
Fu L, Pelicano H, Liu J, Huang P, Lee CC. The circadian gene Period2 plays an important role in tumor suppression and DNA damage response in vivo. \textit{Cell}. 2002;111(1):41--50.

\bibitem{janich2021}
Janich P, et al. The circadian molecular clock creates epidermal stem cell heterogeneity. \textit{Nature}. 2011;480(7376):209--214.

\bibitem{matsuo2003}
Matsuo T, et al. Control mechanism of the circadian clock for timing of cell division in vivo. \textit{Science}. 2003;302(5643):255--259.

\bibitem{mcgowan1995}
McGowan CH, Russell P. Cell cycle regulation of human WEE1. \textit{EMBO J}. 1995;14(10):2166--2175.

\bibitem{altman2015}
Altman BJ, et al. MYC disrupts the circadian clock and metabolism in cancer cells. \textit{Cell Metab}. 2015;22(6):1009--1019.

\bibitem{leijen2016}
Leijen S, et al. Phase II study of WEE1 inhibitor AZD1775 plus carboplatin in patients with TP53-mutated ovarian cancer refractory or resistant to first-line therapy. \textit{J Clin Oncol}. 2016;34(36):4354--4361.

\bibitem{panda2016}
Panda S. Circadian physiology of metabolism. \textit{Science}. 2016;354(6315):1008--1015.

\bibitem{schernhammer2013}
Schernhammer ES, et al. Night-shift work and risk of colorectal cancer in the nurses' health study. \textit{J Natl Cancer Inst}. 2003;95(11):825--828.

\bibitem{ripperger2000clock}
Ripperger JA, Shearman LP, Reppert SM, Schibler U. CLOCK, an essential pacemaker component, controls expression of the circadian transcription factor DBP. \textit{Genes Dev}. 2000;14(6):679--689.

\bibitem{wuarin1990dbp}
Wuarin J, Schibler U. Expression of the liver-enriched transcriptional activator protein DBP follows a stringent circadian rhythm. \textit{Cell}. 1990;63(6):1257--1266.

\bibitem{gachon2006parbzip}
Gachon F, Olela FF, Schaad O, Descombes P, Schibler U. The circadian PAR-domain basic leucine zipper transcription factors DBP, TEF, and HLF modulate basal and inducible xenobiotic detoxification. \textit{Cell Metab}. 2006;4(1):25--36.

\bibitem{reick2001npas2}
Reick M, Garcia JA, Dudley C, McKnight SL. NPAS2: an analog of clock operative in the mammalian forebrain. \textit{Science}. 2001;293(5529):506--509.

\bibitem{debruyne2007clock}
DeBruyne JP, Weaver DR, Reppert SM. CLOCK and NPAS2 have overlapping roles in the suprachiasmatic circadian clock. \textit{Nat Neurosci}. 2007;10(5):543--545.

\bibitem{takeda2012rorc}
Takeda Y, Jothi R, Birault V, Jetten AM. ROR$\gamma$ directly regulates the circadian expression of clock genes and downstream targets in vivo. \textit{Nucleic Acids Res}. 2012;40(17):8519--8535.

\bibitem{siu2010cycline}
Siu KT, Rosner MR, Minella AC. An integrated view of cyclin E function and regulation. \textit{Cell Div}. 2010;5:2.

\bibitem{farshadi2020clockcellcycle}
Farshadi E, van der Horst GTJ, Chaves I. Molecular links between the circadian clock and the cell cycle. \textit{J Mol Biol}. 2020;432(12):3515--3524.

\end{thebibliography}

\newpage
\section*{Supplementary Tables}

\subsection*{Table S1: Complete tissue-by-tissue significant findings}

\begin{longtable}{lllcc}
\toprule
\textbf{Tissue} & \textbf{Clock Gene} & \textbf{Target Gene} & \textbf{Conservation} & \textbf{P-value} \\
\midrule
\endfirsthead
\multicolumn{5}{c}{\textit{Table S1 continued}} \\
\toprule
\textbf{Tissue} & \textbf{Clock Gene} & \textbf{Target Gene} & \textbf{Conservation} & \textbf{P-value} \\
\midrule
\endhead
\midrule
\multicolumn{5}{r}{\textit{Continued on next page}} \\
\endfoot
\bottomrule
\endlastfoot
Liver & Cry1 & Wee1 & Conserved ($\geq$6) & 0.0104 \\
Liver & Nr1d1 & Wee1 & Moderate (3-5) & 0.0104 \\
Liver & Cry2 & Wee1 & Moderate (3-5) & 0.0110 \\
Liver & Per2 & Ccnd1 & Tissue-specific & 0.0150 \\
Brown Fat & Cry1 & Wee1 & Conserved ($\geq$6) & 0.0100 \\
Brown Fat & Nr1d2 & Chek2 & Tissue-specific & 0.0076 \\
White Fat & Cry1 & Wee1 & Conserved ($\geq$6) & 0.0130 \\
White Fat & Clock & Bax & Tissue-specific & 0.0059 \\
Aorta & Cry1 & Wee1 & Conserved ($\geq$6) & 0.0105 \\
Aorta & Nr1d1 & Wee1 & Moderate (3-5) & 0.0038 \\
Heart & Cry1 & Wee1 & Conserved ($\geq$6) & 0.0095 \\
Heart & Nr1d1 & Tead1 & Tissue-specific & 0.0061 \\
Lung & Cry1 & Wee1 & Conserved ($\geq$6) & 0.0190 \\
Lung & Arntl & Wee1 & Moderate (3-5) & 0.0178 \\
Brainstem & Cry1 & Wee1 & Conserved ($\geq$6) & 0.0085 \\
Brainstem & Cry1 & Tead1 & Tissue-specific & 0.0048 \\
Muscle & Cry1 & Wee1 & Conserved ($\geq$6) & 0.0310 \\
Muscle & Nr1d2 & Myc & Tissue-specific & 0.0237 \\
Adrenal & Cry1 & Wee1 & Conserved ($\geq$6) & 0.0015 \\
Adrenal & Nr1d2 & Wee1 & Moderate (3-5) & 0.0006 \\
Cerebellum & Cry1 & Wee1 & Conserved ($\geq$6) & 0.0200 \\
Cerebellum & Nr1d1 & Cdk1 & Tissue-specific & 0.0110 \\
Hypothalamus & Cry1 & Wee1 & Conserved ($\geq$6) & 0.0480 \\
Hypothalamus & Arntl & Chek2 & Tissue-specific & 0.0436 \\
Kidney & Nr1d2 & Myc & Tissue-specific & 0.0431 \\
\end{longtable}

\subsection*{Table S2: Conservation classification summary}

\begin{table}[H]
\centering
\begin{tabular}{lcc}
\toprule
\textbf{Conservation Level} & \textbf{Unique Pairs} & \textbf{Examples} \\
\midrule
Conserved ($\geq$6 tissues) & 1 & Cry1$\rightarrow$Wee1 \\
Moderate (3--5 tissues) & 8 & Nr1d1$\rightarrow$Wee1, Arntl$\rightarrow$Wee1 \\
Tissue-specific (1--2 tissues) & 56 & All Myc, Tead1, Chek2 interactions \\
\bottomrule
\end{tabular}
\end{table}

\end{document}
